% Generated draft: requires mathematical and source review.
\documentclass[11pt]{article}
\usepackage[margin=25mm]{geometry}
\usepackage{fontspec}
\setmainfont{Latin Modern Roman}
\usepackage{amsmath,amssymb,mathtools,unicode-math}
\setmathfont{Latin Modern Math}
\usepackage{enumitem,needspace,xurl,hyperref}
\setlength{\parindent}{0pt}
\setlength{\parskip}{5pt}
\setlength{\emergencystretch}{3em}
\newcommand{\N}{\mathbb{N}}
\newcommand{\Z}{\mathbb{Z}}
\newcommand{\Q}{\mathbb{Q}}
\newcommand{\R}{\mathbb{R}}
\newcommand{\C}{\mathbb{C}}
\newcommand{\F}{\mathbb{F}}
\newcommand{\A}{\mathbb{A}}
\newcommand{\PP}{\mathbb{P}}
\newcommand{\E}{\mathbb{E}}
\newcommand{\eps}{\varepsilon}
\newcommand{\dd}{\,\mathrm{d}}
\begin{document}
\section*{UnsolvedMath ID: 30005384}
\textbf{Source problem code:} catalogue entry
\subsection{Definitions and mathematical statement}
Let \(M\) and \(N\) be compact, connected, oriented topological manifolds of dimension \(4\), as assumed in the supplied source. A homeomorphism \(h\colon M\to N\) is a continuous bijection with continuous inverse. Thus the ordinary homeomorphism relation is
\[
M\cong_{\mathrm{Top}}N
\quad\Longleftrightarrow\quad
\exists h\colon M\to N\text{ such that }h\text{ is a homeomorphism}.
\]
The extracted source asks how one can try to determine whether this relation holds for two such manifolds. It does not specify whether the homeomorphism must preserve the given orientations. If preservation is intended, the existential condition must additionally require that \(h\) preserve orientation. The source also does not specify whether manifolds with boundary are allowed or whether any boundary identifications must be respected.

\textbf{Unresolved target.} The supplied statement-verification note identifies the extracted question as an expository heading immediately followed by a modified-surgery strategy, rather than a standalone open classification problem. The registry marks the statement as \texttt{unrecoverable} and supplies no clean replacement. Consequently, the relation above defines the subject of the discussion; it is not a recovered conjecture or a claim that a general decision algorithm is requested or exists. The intended classification target, allowable input descriptions, and required form of a decision procedure remain unspecified.
\subsection{Short English statement}
The source discusses how to determine whether two compact, connected, oriented topological four-manifolds are homeomorphic. Its extracted question is identified by the supplied verification note as an expository heading, so no precise standalone open problem is recoverable. Orientation preservation, boundary conventions, and any algorithmic requirements remain unresolved.
\subsection{Sources}
\begin{itemize}
\item Supplied local UnsolvedMath registry record, permanent ID \(30005384\), source problem code \texttt{catalogue entry}, titled ``Homeomorphism Classification of Compact Oriented Four-Manifolds.'' The record belongs to the Oberwolfach Reports Open Problems collection and gives proposed year \(2023\); no individual proposer is supplied.
\item \emph{Morphisms in Low Dimensions} (2023), Oberwolfach Reports, Report \(4/2023\), pp.\ 215--259, DOI \texttt{10.4171/owr/2023/4}: \url{https://doi.org/10.4171/owr/2023/4}. This is the supplied report-level citation; no precise location for the extracted heading is provided.
\item Conway, Crowley, Powell, and Sixt, \emph{Large homotopy stable classes}: \url{https://arxiv.org/abs/2109.00654}. This is supplied as a literature-review reference, not as an identified source of a standalone problem statement.
\item Literature assessment checked \texttt{2026-08-22}: supplied dated curation, not independently verified. The record labels the topic ``open'' and reports that stable surgery and homotopy equivalence provide useful decision procedures in restricted settings, while no general classification or homeomorphism test for arbitrary compact oriented four-manifolds is known. The background explicitly requires independent verification of time-sensitive status and cited claims.
\item The supplied statement-verification note says that the extracted sentence is an expository heading immediately followed by the modified-surgery strategy, not a standalone open classification problem. It records that no separate clean statement is promoted and that the original record remains publicly visible for provenance.
\end{itemize}
\end{document}
